\chapter{Calcul Littéral \& Expressions}
\label{chap:calcul_litteral}

% ==============================================================================
% 1. OBJECTIFS & COMPÉTENCES DU RÉFÉRENTIEL
% ==============================================================================
\begin{cadreretenu}[Objectifs \& Compétences du DNB Pro]
\begin{itemize}
    \item \textbf{Connaissances} : Conventions d'écriture littérale, vocabulaire (somme, produit, terme, facteur), identité remarquable $a^2 - b^2 = (a-b)(a+b)$.
    \item \textbf{Capacités} :
    \begin{itemize}
        \item \capRealiser\ Réduire une expression algébrique en regroupant les termes de même nature.
        \item \capRealiser\ Développer une expression en utilisant la distributivité simple $k(a+b)$ et double $(a+b)(c+d)$.
        \item \capRealiser\ Factoriser une expression simple (facteur commun évident ou différence de deux carrés).
        \item \capAnalyser\ Traduire une situation concrète (périmètre, aire, coût de production) par une formule littérale.
        \item \capValider\ Tester l'égalité de deux expressions en substituant des valeurs numériques.
    \end{itemize}
\end{itemize}
\end{cadreretenu}

\vspace{0.3cm}

% ==============================================================================
% 2. ACTIVITÉ D'INVESTIGATION
% ==============================================================================
\section{Activité d'Introduction : L'Atelier de Menuiserie}

\begin{activite}{Conception d'une structure en bois modulable}
Un menuisier conçoit des cadres en bois rectangulaires dont la largeur est notée $x$ (en mètres) et la longueur mesure $3\text{ m}$ de plus que la largeur.

\begin{center}
\begin{tikzpicture}[scale=0.9]
    \draw[fill=amber!20, draw=amber!80!black, line width=1.2pt] (0,0) rectangle (6,3.5);
    \draw[<->, >=stealth, slate700, thick] (0,-0.4) -- (6,-0.4) node[midway, below] {Longueur : $L = x + 3$};
    \draw[<->, >=stealth, slate700, thick] (-0.4,0) -- (-0.4,3.5) node[midway, left] {Largeur : $\ell = x$};
    \node at (3,1.75) {\textbf{Cadre en Chêne}};
\end{tikzpicture}
\end{center}

\begin{enumerate}
    \item \capSApproprier\ Exprimer le périmètre $P(x)$ du cadre en fonction de la variable $x$ sous forme d'une somme :
    \[ P(x) = x + (x + 3) + x + (x + 3) \]
    \item \capRealiser\ Réduire cette expression en regroupant les termes en $x$ et les nombres :
    \lignesreponse[2]
    \item \capRealiser\ Exprimer l'aire $A(x)$ du rectangle en fonction de $x$ sous la forme d'un produit :
    \[ A(x) = x \times (x + 3) \]
    \item \capAnalyser\ Développer ce produit pour obtenir l'expression de l'aire sous forme d'une somme :
    \lignesreponse[2]
    \item \capValider\ Calculer le périmètre et l'aire pour une largeur $x = 2\text{ m}$ :
    \lignesreponse[2]
\end{enumerate}
\end{activite}

% ==============================================================================
% 3. COURS ET NOTIONS ESSENTIELLES
% ==============================================================================
\section{Cours : Conventions, Développement et Réduction}

\subsection{Conventions d'écriture et Réduction}

\begin{cadredef}[Convention de simplification]
Dans une expression algébrique, le signe de multiplication $\times$ peut être omis devant une lettre ou une parenthèse :
\[ 3 \times x = 3x \qquad 1 \times x = x \qquad x \times x = x^2 \qquad 5 \times (x + 2) = 5(x + 2) \]
\end{cadredef}

\begin{cadremethode}[Réduire une somme littérale]
Pour réduire une expression littérale, on regroupe et on additionne les termes de \textbf{même nature} (les termes en $x^2$ ensemble, les termes en $x$ ensemble, les constantes numériques ensemble).
\end{cadremethode}

\begin{cadreexemple}[Réduction guidée]
Réduire l'expression $A = 5x^2 + 3x - 7 - 2x^2 + 8x + 4$ :
\begin{align*}
    A &= (5x^2 - 2x^2) + (3x + 8x) + (-7 + 4) \\
    A &= 3x^2 + 11x - 3
\end{align*}
\end{cadreexemple}

\begin{cadrepiege}[Erreur classique d'addition]
On ne peut \textbf{jamais} additionner des termes de familles différentes !
\[ 3x + 4 \neq 7x \qquad\text{et}\qquad 2x^2 + 3x \neq 5x^3 \]
\end{cadrepiege}

\subsection{Développement : Distributivité Simple et Double}

\begin{cadreprop}[Formules de distributivité]
Pour tous nombres réels $k$, $a$, $b$, $c$ et $d$ :
\begin{itemize}
    \item \textbf{Distributivité Simple} :
    \[ k(a + b) = k \times a + k \times b = ka + kb \]
    \[ k(a - b) = k \times a - k \times b = ka - kb \]
    \item \textbf{Double Distributivité} :
    \[ (a + b)(c + d) = ac + ad + bc + bd \]
\end{itemize}
\end{cadreprop}

\begin{cadreexemple}[Développement d'une double distributivité]
Développer et réduire l'expression $B = (2x + 3)(x - 5)$ :
\begin{align*}
    B &= 2x \times x + 2x \times (-5) + 3 \times x + 3 \times (-5) \\
    B &= 2x^2 - 10x + 3x - 15 \\
    B &= 2x^2 - 7x - 15
\end{align*}
\end{cadreexemple}

\subsection{Factorisation et Identité Remarquable}

\begin{cadredef}[Factoriser]
\textbf{Factoriser} une expression consiste à transformer une somme (ou différence) en un \textbf{produit de facteurs}.
\[ ka + kb = k(a + b) \]
\end{cadredef}

\begin{cadreprop}[Identité remarquable : Différence de deux carrés]
Pour tous nombres réels $a$ et $b$ :
\[ a^2 - b^2 = (a - b)(a + b) \]
\end{cadreprop}

\begin{cadreexemple}[Factorisations usuelles]
\begin{itemize}
    \item Par facteur commun : $C = 6x + 18 = 6 \times x + 6 \times 3 = 6(x + 3)$
    \item Par identité remarquable : $D = x^2 - 49 = x^2 - 7^2 = (x - 7)(x + 7)$
\end{itemize}
\end{cadreexemple}

% ==============================================================================
% 4. AUTOMATISMES & QUESTIONS FLASH
% ==============================================================================
\section{Questions Flash / Automatismes}

\begin{automatisme}[Calcul Mental \& Réflexes Algébriques (10 min)]
\noindent
\begin{tabularx}{\linewidth}{X|X}
\textbf{1.} Réduire : $4x + 7x = \pointille[2cm]$ & \textbf{6.} Développer : $3(2x + 5) = \pointille[2.5cm]$ \\
\textbf{2.} Réduire : $5x^2 - 9x^2 = \pointille[2cm]$ & \textbf{7.} Développer : $-2(x - 4) = \pointille[2.5cm]$ \\
\textbf{3.} Calculer $A = 3x - 2$ pour $x = 4$ : & \textbf{8.} Factoriser : $7x + 21 = \pointille[2.5cm]$ \\
$\implies A = \pointille[2cm]$ & \textbf{9.} Factoriser : $x^2 - 16 = \pointille[2.5cm]$ \\
\textbf{4.} Calculer $B = x^2 + 5$ pour $x = -3$ : & \textbf{10.} Vrai ou Faux ? \\
$\implies B = \pointille[2cm]$ & $(x+2)^2 = x^2 + 4$ : \pointille[2cm] \\
\textbf{5.} Simplifier : $3x \times 4x = \pointille[2cm]$ & \textbf{Bonus :} Développer $(x+1)(x+3) = \pointille[2.5cm]$ \\
\end{tabularx}
\end{automatisme}

% ==============================================================================
% 5. TRAVAUX DIRIGÉS (TD) - 10 EXERCICES CORRIGÉS
% ==============================================================================
\section{Travaux Dirigés (TD) : Exercices d'Entraînement}

\begin{exercice}[2 pts]{\capRealiser}{Simplification et Réduction}
Réduire les expressions suivantes au maximum :
\begin{enumerate}
    \item $A = 7x + 3 - 2x + 9$
    \item $B = 4x^2 - 6x + 8 - x^2 + 10x - 12$
\end{enumerate}
\lignesreponse[4]
\end{exercice}

\begin{exercice}[2 pts]{\capRealiser}{Développement par Distributivité Simple}
Développer puis réduire les expressions :
\begin{enumerate}
    \item $C = 4(3x - 5)$
    \item $D = -3x(2x - 7)$
    \item $E = 5(x + 2) - 2(3x - 4)$
\end{enumerate}
\lignesreponse[6]
\end{exercice}

\begin{exercice}[3 pts]{\capRealiser}{Double Distributivité}
Développer, réduire et ordonner suivant les puissances décroissantes de $x$ :
\begin{enumerate}
    \item $F = (x + 4)(x + 7)$
    \item $G = (2x - 3)(x + 5)$
    \item $H = (3x - 2)(4x - 1)$
\end{enumerate}
\lignesreponse[6]
\end{exercice}

\begin{exercice}[2 pts]{\capRealiser}{Factorisation et Identités}
Factoriser chacune des expressions suivantes :
\begin{enumerate}
    \item $I = 5x + 35$
    \item $J = x^2 - 81$
    \item $K = 4x^2 - 25$
\end{enumerate}
\lignesreponse[5]
\end{exercice}

\begin{exercice}[3 pts]{\capAnalyser}{Problème Métier BTP : Périmètre d'une clôture}
Sur un chantier de construction, un chef d'équipe doit sécuriser une zone rectangulaire de longueur $L = 3x + 10$ mètres et de largeur $\ell = 2x + 4$ mètres.
\begin{enumerate}
    \item Exprimer le périmètre $P$ de la zone en fonction de $x$ sous forme développée et réduite.
    \item Calculer la longueur totale de grillage nécessaire si $x = 6\text{ m}$.
    \item Le grillage est vendu en rouleaux de $25\text{ m}$ à $48\text{ \euro}$ l'unité. Déterminer le coût total d'achat.
\end{enumerate}
\lignesreponse[6]
\end{exercice}

\begin{exercice}[3 pts]{\capAnalyser}{Problème Aménagement : Aire d'une terrasse}
Un paysagiste aménage une terrasse dont l'aire totale est donnée par la formule $A(x) = (x + 5)(2x + 3) - (x^2 + 15)$.
\begin{enumerate}
    \item Développer et réduire l'expression $A(x)$.
    \item Calculer l'aire de la terrasse pour $x = 4\text{ m}$.
\end{enumerate}
\lignesreponse[5]
\end{exercice}

% ==============================================================================
% 6. TP TICE TABLEUR
% ==============================================================================
\section{TP TICE : Modélisation et Tableur}

\begin{tpinfo}[Tableur LibreOffice Calc / Excel]{Formules et Équivalence Algébrique}
Un artisan teste deux tarifs d'installation électrique :
\begin{itemize}
    \item \textbf{Formule A} : $50\text{ \euro}$ de forfait déplacement $+ 35\text{ \euro}$ par heure : $T_A(x) = 35x + 50$
    \item \textbf{Formule B} : $20\text{ \euro}$ de forfait $+ 42\text{ \euro}$ par heure : $T_B(x) = 42x + 20$
\end{itemize}

\begin{center}
\begin{tabular}{|c|c|c|c|}
\hline
\rowcolor{slate200} & \textbf{A} & \textbf{B} & \textbf{C} \\
\hline
\textbf{1} & \textbf{Nombre d'heures ($x$)} & \textbf{Tarif A (\euro)} & \textbf{Tarif B (\euro)} \\
\hline
\textbf{2} & 1 & \texttt{=35*A2+50} & \texttt{=42*A2+20} \\
\hline
\textbf{3} & 2 & & \\
\hline
\textbf{4} & 3 & & \\
\hline
\textbf{5} & 4 & & \\
\hline
\textbf{6} & 5 & & \\
\hline
\end{tabular}
\end{center}

\begin{enumerate}
    \item Quelle formule tableur doit-on saisir dans la cellule \texttt{B2} puis étirer vers le bas ?
    \item À partir de combien d'heures de travail la Formule A devient-elle plus avantageuse pour le client ?
\end{enumerate}
\lignesreponse[3]
\end{tpinfo}

% ==============================================================================
% 7. VERS LE BREVET (DNB PRO)
% ==============================================================================
\section{Vers le Brevet (DNB Pro)}

\begin{sujetdnb}[DNB Pro Métropole][6 pts]{Menuiserie \& Terrasse en Bois}
Un artisan menuisier réalise une terrasse en lames de bois composite. Le plan ci-dessous indique les dimensions exprimées en mètres :

\begin{center}
\begin{tikzpicture}[scale=0.85]
    \draw[fill=colDNBBg, draw=colDNB, line width=1pt] (0,0) -- (7,0) -- (7,3) -- (4,3) -- (4,5) -- (0,5) -- cycle;
    \node at (2,2.5) {\textbf{Partie 1}};
    \node at (5.5,1.5) {\textbf{Partie 2}};
    \draw[<->, >=stealth, slate700] (0,-0.3) -- (7,-0.3) node[midway, below] {$x + 4$};
    \draw[<->, >=stealth, slate700] (-0.3,0) -- (-0.3,5) node[midway, left] {$x + 2$};
    \draw[<->, >=stealth, slate700] (7.3,0) -- (7.3,3) node[midway, right] {$3$};
    \draw[<->, >=stealth, slate700] (4,5.3) -- (0,5.3) node[midway, above] {$4$};
\end{tikzpicture}
\end{center}

\begin{enumerate}
    \item \capSApproprier\ Montrer que l'aire totale de la terrasse s'exprime par la relation :
    \[ \mathcal{A}(x) = 4(x + 2) + 3x \]
    \item \capRealiser\ Développer et réduire cette expression.
    \item \capValider\ Pour une terrasse réelle avec $x = 3\text{ m}$ :
    \begin{enumerate}
        \item Calculer l'aire exacte en $\text{m}^2$.
        \item Les lames coûtent $32\text{ \euro}/\text{m}^2$. Calculer le montant total du bois nécessaire.
    \end{enumerate}
\end{enumerate}
\lignesreponse[6]
\end{sujetdnb}

% ==============================================================================
% 8. CORRIGÉS DÉTAILLÉS
% ==============================================================================
\section{Corrigés Détaillés du Chapitre}

\begin{solution}[Corrigé des Exercices de TD]
\begin{itemize}
    \item \textbf{Exercice 1} :
    \begin{enumerate}
        \item $A = (7x - 2x) + (3 + 9) = \repbox{5x + 12}$
        \item $B = (4x^2 - x^2) + (-6x + 10x) + (8 - 12) = \repbox{3x^2 + 4x - 4}$
    \end{enumerate}
    \item \textbf{Exercice 2} :
    \begin{enumerate}
        \item $C = 4 \times 3x + 4 \times (-5) = \repbox{12x - 20}$
        \item $D = -3x \times 2x - 3x \times (-7) = \repbox{-6x^2 + 21x}$
        \item $E = 5x + 10 - 6x + 8 = \repbox{-x + 18}$
    \end{enumerate}
    \item \textbf{Exercice 3} :
    \begin{enumerate}
        \item $F = x^2 + 7x + 4x + 28 = \repbox{x^2 + 11x + 28}$
        \item $G = 2x^2 + 10x - 3x - 15 = \repbox{2x^2 + 7x - 15}$
        \item $H = 12x^2 - 3x - 8x + 2 = \repbox{12x^2 - 11x + 2}$
    \end{enumerate}
    \item \textbf{Exercice 4} :
    \begin{enumerate}
        \item $I = 5(x + 7)$
        \item $J = x^2 - 9^2 = (x - 9)(x + 9)$
        \item $K = (2x)^2 - 5^2 = (2x - 5)(2x + 5)$
    \end{enumerate}
    \item \textbf{Exercice 5 (BTP)} :
    \begin{enumerate}
        \item $P = 2(L + \ell) = 2((3x + 10) + (2x + 4)) = 2(5x + 14) = \repbox{10x + 28}\text{ m}$.
        \item Pour $x = 6$ : $P = 10 \times 6 + 28 = 60 + 28 = \repbox{88\text{ m}}$.
        \item Nombre de rouleaux : $88 / 25 = 3{,}52 \implies$ il faut \textbf{4 rouleaux}. Coût : $4 \times 48 = \repbox{192\text{ \euro}}$.
    \end{enumerate}
\end{itemize}
\end{solution}

\begin{solution}[Corrigé du Sujet DNB Pro]
\begin{enumerate}
    \item L'aire est la somme de l'aire du rectangle 1 ($4 \times (x + 2)$) et du rectangle 2 ($3 \times x$). Donc $\mathcal{A}(x) = 4(x+2) + 3x$.
    \item $\mathcal{A}(x) = 4x + 8 + 3x = \repbox{7x + 8}$.
    \item Pour $x = 3\text{ m}$ :
    \begin{enumerate}
        \item $\mathcal{A}(3) = 7 \times 3 + 8 = 21 + 8 = \repbox{29\text{ m}^2}$.
        \item Coût du bois : $29 \times 32 = \repbox{928\text{ \euro}}$.
    \end{enumerate}
\end{enumerate}
\end{solution}
